\chapter{Search and Unemployment}\label{ch:b2:search-and-unemployment}

At any moment some workers are looking for jobs while some firms are looking for workers. The Year 1 volume described the
flows between employment and unemployment; this chapter builds a model in which both sides search, wages are bargained,
and the unemployment rate is an equilibrium outcome. The chart shows the two sides of the euro area's labour market,
quarter by quarter, since 2009.

\begin{openingchart}{Euro area: job vacancy rate (industry, construction and services) against the unemployment rate (ages
  15--74), quarterly, seasonally adjusted, 2009--2025.}{Eurostat (jvs\_q\_nace2, une\_rt\_q); our computation}
\begin{tikzpicture}
\begin{axis}[omrecord, width=0.62\linewidth, height=0.46\linewidth, xmin=5, xmax=13.5, ymin=0.5, ymax=3.6,
    xlabel={Unemployment rate, per cent}, ylabel={Job vacancy rate, per cent},
    legend style={at={(0.97,0.97)}, anchor=north east, font=\scriptsize}]
  \addplot[omDef, thick, mark=*, mark size=1.1pt, restrict expr to domain={\thisrow{period}}{1:1}] table[x=u, y=v]
    {figdata/bachelor-2/26-search-and-unemployment/out/beveridge.csv};
  \addlegendentry{2009--2019}
  \addplot[omThm, thick, mark=square*, mark size=1.1pt, restrict expr to domain={\thisrow{period}}{0:0}] table[x=u, y=v]
    {figdata/bachelor-2/26-search-and-unemployment/out/beveridge.csv};
  \addlegendentry{2020--2025}
  % positions of the labelled quarters: out/bev_points.csv
  \node[font=\scriptsize, anchor=north west] at (axis cs:9.0,1.0) {2009-Q1};
  \node[font=\scriptsize, anchor=west] at (axis cs:12.35,1.3) {2013-Q1};
  \node[font=\scriptsize, anchor=north] at (axis cs:7.4,2.27) {2019-Q4};
  \node[font=\scriptsize, anchor=north west] at (axis cs:8.7,1.8) {2020-Q3};
  \node[font=\scriptsize, anchor=south] at (axis cs:6.7,3.32) {2022-Q2};
  \node[font=\scriptsize, anchor=east] at (axis cs:6.2,2.2) {2025-Q4};
\end{axis}
\end{tikzpicture}
\end{openingchart}

% figdata: bev_summary (2009-Q1 9.0/1.0; 2013-Q1 12.3/1.3; 2019-Q4 7.4/2.3, theta 0.311; 2020-Q3 8.7/1.8; 2022-Q2 6.7/3.3, theta 0.493; 2025-Q4 6.3/2.2)
From 2013 to 2019 unemployment fell from 12.3 to 7.4 per cent while the vacancy rate rose from 1.3 to 2.3 per cent: the economy
moved up and to the left along a downward-sloping curve. In 2020 vacancies fell and unemployment rose to 8.7 per cent; then
vacancies rose to 3.3 per cent in the second quarter of 2022, far above anything seen before, with unemployment at 6.7. By the end of
2025 the vacancy rate was back at 2.2 per cent with unemployment still lower, at 6.3. The ratio of vacancies to unemployment went
from 0.31 at the end of 2019 to 0.49 in mid-2022.

\begin{recall}
From the Year 1 volume: the separation rate and the job-finding rate, the steady-state unemployment rate, the Beveridge
curve and efficiency wages. From \cref{ch:b2:expectations-and-lucas-critique}: equilibrium with forward-looking agents.
\end{recall}

\begin{notation}[Symbols of this chapter]\label{not:b2:search-and-unemployment:symbols}
Time is continuous; rates are per period. $u$ is the unemployment rate and $v$ the vacancy rate, both as shares of the labour force; $s$ is the separation rate; $r$ the interest rate;
$p$ output per worker; $z$ the value of unemployment per period (benefits and leisure, a local symbol here); $c$ the cost of an open vacancy per period; $\eta$ the
elasticity of matches with respect to unemployment; $\beta$ the worker's \omterm{def:b2:search-and-unemployment:nash-bargaining}{bargaining power} (a local symbol here, not the
\omterm{def:b2:intertemporal-choice:discount-factor}{discount factor}).
\end{notation}

\section{The matching function}\label{sec:b2:search-and-unemployment:matching}

\begin{definition}[Matching function]\label{def:b2:search-and-unemployment:matching-function}
A \emph{matching function}\index{matching function} gives the number of new hires per period as a function of the number of
unemployed workers and the number of vacancies: $m = m(u, v)$, increasing in both.
\end{definition}

The \omterm{def:b2:search-and-unemployment:matching-function}{matching function} summarises frictions the model does not describe: workers and firms do not know where the others are,
skills and places differ, and applications take time. The usual form has constant returns, $m = Au^\eta v^{1-\eta}$, so that doubling both sides doubles
hires.

\begin{definition}[Labour-market tightness]\label{def:b2:search-and-unemployment:tightness}
\emph{Labour-market tightness}\index{labour-market tightness} is the ratio of vacancies to unemployed workers, $\theta = v/u$.
\end{definition}

\begin{definition}[Vacancy-filling rate]\label{def:b2:search-and-unemployment:vacancy-filling}
The \emph{vacancy-filling rate}\index{vacancy-filling rate} is the probability per period that an open vacancy is filled,
$q(\theta) = m/v = A\theta^{-\eta}$.
\end{definition}

With constant returns everything depends on tightness. The job-finding rate is $f(\theta) = m/u = A\theta^{1-\eta} = \theta q(\theta)$: in a tight market workers find jobs fast and firms fill vacancies
slowly. With $A = 0.7$, $u = 0.06$ and $v = 0.03$, $\theta = 0.5$, matches are $0.7 \times \sqrt{0.06 \times 0.03} = 0.0297$, the job-finding rate is 0.495 and the \omterm{def:b2:search-and-unemployment:vacancy-filling}{vacancy-filling rate} 0.990.

\section{Flows and the Beveridge curve}\label{sec:b2:search-and-unemployment:beveridge}

Unemployment changes as $\dot u = s(1 - u) - f(\theta)u$. In a steady state inflows equal outflows:
\begin{equation}\label{eq:b2:search-and-unemployment:bc}
u = \frac{s}{s + f(\theta)}.
\end{equation}
Given $s$, each tightness gives one unemployment rate, and with $v = \theta u$ one vacancy rate. As $\theta$ rises, $u$ falls and $v$ rises: the steady states
trace a downward-sloping curve in $(u, v)$, the Beveridge curve of the opening chart. A fall in matching efficiency $A$ or a rise in $s$ shifts it
outward: more of both at once.

\section{Job creation}\label{sec:b2:search-and-unemployment:job-creation}

A firm that opens a vacancy pays $c$ per period until it is filled, which takes $1/q(\theta)$ periods on average. A filled job produces $p$ and pays $w$ until it
is destroyed at rate $s$; its value, discounted at $r$, is $(p - w)/(r + s)$. Firms open vacancies until the expected value of one more is zero.

\begin{definition}[Job-creation condition]\label{def:b2:search-and-unemployment:job-creation}
The \emph{job-creation condition}\index{job-creation condition} states that, with free entry of vacancies, the expected cost of filling a vacancy equals the value of a filled job:
$\dfrac{c}{q(\theta)} = \dfrac{p - w}{r + s}$.
\end{definition}

A higher wage makes jobs less valuable, so firms open fewer vacancies until the faster filling of each restores the
balance: the job-creation curve slopes down in $(\theta, w)$.

\section{Wage bargaining}\label{sec:b2:search-and-unemployment:bargaining}

Once a worker and a firm meet, a match produces more than the two would get apart; both would lose by walking away and
searching again. The wage divides this surplus.

\begin{definition}[Nash bargaining solution]\label{def:b2:search-and-unemployment:nash-bargaining}
The \emph{Nash bargaining solution}\index{Nash bargaining solution} of a bargain over a surplus sets the terms that
maximise $(W - U)^\beta(J - V)^{1-\beta}$, where $W - U$ is the worker's gain from the match and $J - V$ the firm's; $\beta \in (0, 1)$ is the worker's \emph{bargaining power}\index{bargaining power},
and the solution gives the worker the share $\beta$ of the total surplus.
\end{definition}

With free entry, $V = 0$, and with the value equations of the worker and the firm (Exercise 11) the Nash solution gives
\begin{equation}\label{eq:b2:search-and-unemployment:wage}
w = (1 - \beta)z + \beta(p + c\theta).
\end{equation}
The wage is a weighted average of what the worker gets unemployed and what the match produces, plus the hiring costs
the firm saves by keeping the worker, which are larger in a tight market. This wage curve slopes up in $(\theta, w)$.

\section{Equilibrium unemployment}\label{sec:b2:search-and-unemployment:dmp}

\begin{definition}[Diamond--Mortensen--Pissarides model]\label{def:b2:search-and-unemployment:dmp}
The \emph{Diamond--Mortensen--Pissarides model}\index{Diamond--Mortensen--Pissarides model} is the equilibrium model of
unemployment that combines a \omterm{def:b2:search-and-unemployment:matching-function}{matching function}, free entry of vacancies and wages set by Nash bargaining.
\end{definition}

Substituting the wage into the \omterm{def:b2:search-and-unemployment:job-creation}{job-creation condition} gives one equation in $\theta$:
\begin{equation}\label{eq:b2:search-and-unemployment:theta}
\frac{(1 - \beta)(p - z)}{r + s + \beta f(\theta)} = \frac{c}{q(\theta)}.
\end{equation}
The left side falls with $\theta$ and the right side rises, so there is one equilibrium tightness. It sets the wage through
\eqref{eq:b2:search-and-unemployment:wage} and unemployment through \eqref{eq:b2:search-and-unemployment:bc}.

\begin{omfigure}
\begin{tikzpicture}[econ]
\begin{axis}[omrecord, grid=none, xmin=0, xmax=1.2, ymin=0.78, ymax=1.06, width=0.55\linewidth, height=0.42\linewidth,
    xlabel={Tightness, $\theta$}, ylabel={Wage, $w$}, scaled y ticks=false, yticklabel style={/pgf/number format/fixed},
    legend style={at={(1.03,0.5)}, anchor=west, font=\small}]
  \addplot[omDef, very thick, domain=0:1.2, samples=60] {1 - 0.05*sqrt(x)};
  \addlegendentry{job creation}
  \addplot[omThm, very thick, domain=0:1.04, samples=2] {0.8 + 0.25*x};
  \addlegendentry{wage curve, $z = 0.6$}
  \addplot[omThm, thick, dashed, domain=0:0.84, samples=2] {0.85 + 0.25*x};
  \addlegendentry{wage curve, $z = 0.7$}
  \node[eqpoint] at (axis cs:0.64,0.96) {};
  \node[eqpoint] at (axis cs:0.4638,0.966) {};
  \draw[guide] (axis cs:0.64,0.96) -- (axis cs:0.64,0.78);
  \draw[guide] (axis cs:0.4638,0.966) -- (axis cs:0.4638,0.78);
\end{axis}
\end{tikzpicture}

{\small Brenland's labour market (the weekend problem): equilibrium tightness and wage where the job-creation curve meets
the wage curve. A rise in the value of unemployment from 0.6 to 0.7 shifts the wage curve up: tightness falls from 0.64 to 0.46
and the wage rises from 0.960 to 0.966.}
\end{omfigure}

\begin{omfigure}
\begin{tikzpicture}[econ]
\begin{axis}[omrecord, grid=none, xmin=4, xmax=16, ymin=0, ymax=12, width=0.55\linewidth, height=0.42\linewidth,
    xlabel={Unemployment, per cent}, ylabel={Vacancies, per cent},
    legend style={at={(1.03,0.5)}, anchor=west, font=\small}]
  \addplot[omDef, very thick, domain=4.5:16, samples=80] {69.444*(1 - x/100)^2/x};
  \addlegendentry{Beveridge curve}
  \addplot[omThm, very thick, domain=0:16, samples=2] {0.64*x};
  \addlegendentry{$\theta = 0.64$}
  \addplot[omThm, thick, dashed, domain=0:16, samples=2] {0.4638*x};
  \addlegendentry{$\theta = 0.46$}
  \node[eqpoint] at (axis cs:9.43,6.04) {};
  \node[eqpoint] at (axis cs:10.90,5.06) {};
\end{axis}
\end{tikzpicture}

{\small The same economy in $(u, v)$: tightness is a ray from the origin, and the steady state lies where it crosses the Beveridge curve. The benefit
rise turns the ray down, and unemployment rises from 9.43 to 10.90 per cent.}
\end{omfigure}

The listing solves \eqref{eq:b2:search-and-unemployment:theta} for a monthly calibration with invented numbers, and asks two questions: what a rise in the value of unemployment does,
and how much tightness responds to productivity.

\omcode{code/bachelor-2/26-search-and-unemployment/listings/dmp.py}{1}{30}{Equilibrium tightness, unemployment and
wage.}\label{lst:b2:search-and-unemployment:dmp}

\omoutput{code/bachelor-2/26-search-and-unemployment/outputs/dmp.txt}{A rise in benefits raises unemployment; with $z = 0.4$ a 1 per cent rise in
productivity raises tightness by only 1.8 per cent.}\label{out:b2:search-and-unemployment:dmp}

\begin{model}{Diamond--Mortensen--Pissarides}
\assumes Random matching with constant returns; free entry of vacancies; wages renegotiated by Nash bargaining at every
moment; exogenous separations.
\predicts A higher value of unemployment, through benefits, lowers tightness and raises unemployment and wages; higher
productivity raises tightness and lowers unemployment; vacancies and unemployment move in opposite directions along a
Beveridge curve.
\tested The sign of the benefit effect agrees with natural experiments in which longer potential benefits lowered
job-finding rates (Card, Chetty and Weber 2007; Lalive 2008), though these mainly measure how the unemployed search; the
downward-sloping Beveridge curve of the opening chart.
\breaks Fluctuations: the vacancy-unemployment ratio varies almost 20 times as much as productivity in the American data, while the
model with standard parameters predicts nearly the same volatility for the two (Shimer 2005).
\end{model}

\section{Benefits and productivity}\label{sec:b2:search-and-unemployment:benefits}

% code: listings/dmp.py
In the listing, a rise in $z$ from 0.4 to 0.5 lowers tightness from 0.529 to 0.436 and raises unemployment from 5.56 to 6.10 per cent: a
higher value of unemployment strengthens the worker's position in the bargain, raises wages, and makes vacancies less
profitable. A rise in productivity does the reverse. Both predictions concern the effect of benefits through wages and
vacancies; the evidence below concerns the more direct effect on how hard the unemployed search, which the model's extensions
include.

\begin{record}[How benefits change job finding]
% ledger: b2:ccw2007, b2:lalive2008
Two studies used sharp rules in Austria's unemployment insurance. Card, Chetty and Weber (2007), with data on over half a
million job losers, found that extending the potential duration of benefits from 20 to 30 weeks lowered job-finding rates in the
first 20 weeks of search by 5 to 9 per cent, and that a severance payment of two months' earnings lowered them by 8 to 12 per cent, with
little or no effect of the longer search on the quality of the next job. Lalive (2008) studied a programme that extended
the maximum duration from 30 to 209 weeks for job seekers aged 50 or more in some regions: each extra week of benefits
lengthened unemployment by at least 0.09 weeks for men and 0.32 weeks for women. How large these effects are in other countries
and periods is a question for the Year 3 volume.
\end{record}

\section{Where it breaks: unemployment moves too little}\label{sec:b2:search-and-unemployment:puzzle}

\begin{definition}[Unemployment-volatility puzzle]\label{def:b2:search-and-unemployment:volatility-puzzle}
The \emph{unemployment-volatility puzzle}\index{unemployment-volatility puzzle} is the failure of the
\omterm{def:b2:search-and-unemployment:dmp}{Diamond--Mortensen--Pissarides model}, with standard parameters, to generate fluctuations in tightness and unemployment as
large as those observed, given the observed fluctuations in productivity.
\end{definition}

% ledger: b2:shimer2005; code: listings/dmp.py
In the listing, with $z = 0.4$, a 1 per cent rise in productivity raises tightness by 1.8 per cent. In the American data the
vacancy-unemployment ratio varies almost 20 times as much as productivity (Shimer 2005). The reason is in the wage equation:
when productivity rises, the bargained wage absorbs most of the rise, and the firm's profit from a match, which is what
drives vacancy creation, barely moves. If $z$ is close to $p$, the profit $p - w$ is small, and a 1 per cent rise in $p$ is a large proportional rise in profit: with $z = 0.9$
tightness rises by 11.6 per cent. But that calibration says that a job is worth little more to a worker than unemployment, and in
the listing it puts unemployment at 13.64 per cent unless the vacancy cost is changed too. Wages that respond less to market
conditions than Nash bargaining implies are the other route the literature has taken; neither has settled the question.

\section{Exercises}\label{sec:b2:search-and-unemployment:exercises}

\begin{exercise}[$\star$]\label{exo:b2:search-and-unemployment:1}
Read the opening chart. Give unemployment, vacancies and tightness in 2013-Q1, 2019-Q4 and 2022-Q2.
\end{exercise}

\begin{exercise}[$\star$]\label{exo:b2:search-and-unemployment:2}
With a separation rate of 0.03 and a job-finding rate of 0.5 per month, compute the steady-state unemployment rate.
\end{exercise}

\begin{exercise}[$\star$]\label{exo:b2:search-and-unemployment:3}
With $m = 0.7u^{0.5}v^{0.5}$, $u = 0.06$ and $v = 0.03$, compute matches, tightness, the job-finding rate and the \omterm{def:b2:search-and-unemployment:vacancy-filling}{vacancy-filling rate}.
\end{exercise}

\begin{exercise}[$\star$]\label{exo:b2:search-and-unemployment:4}
Why do firms open vacancies until the expected value of one more is zero, rather than stopping when the last one is
profitable?
\end{exercise}

\begin{exercise}[$\star\star$]\label{exo:b2:search-and-unemployment:5}
With $\beta = 0.5$, $z = 0.4$, $p = 1$, $c = 1$ and $\theta = 0.529$, compute the bargained wage. Which part of it rewards the worker for the firm's hiring cost?
\end{exercise}

\begin{exercise}[$\star\star$]\label{exo:b2:search-and-unemployment:6}
Predict, then check with the listing: the value of unemployment rises from 0.4 to 0.5. Which way do tightness, the wage and unemployment
move?
\end{exercise}

\begin{exercise}[$\star\star$]\label{exo:b2:search-and-unemployment:7}
In Lalive's study, the maximum duration rose by how many weeks? By at least how many weeks did unemployment lengthen, for men
and for women?
\end{exercise}

\begin{exercise}[$\star\star$]\label{exo:b2:search-and-unemployment:8}
Modify the listing to give workers \omterm{def:b2:search-and-unemployment:nash-bargaining}{bargaining power} 0.7. Report tightness, unemployment and the wage, and explain the change.
\end{exercise}

\begin{exercise}[$\star\star$]\label{exo:b2:search-and-unemployment:9}
In the opening chart, compare the fourth quarters of 2019 and 2025. Has the economy moved along a fixed Beveridge curve, or
has the curve shifted? Which parameters of the model would shift it this way, and in which direction?
\end{exercise}

\begin{exercise}[$\star\star\star$]\label{exo:b2:search-and-unemployment:10}
From \eqref{eq:b2:search-and-unemployment:theta} with $\eta = 0.5$, show that $x = \sqrt\theta$ solves a quadratic equation, and write its positive root.
\end{exercise}

\begin{exercise}[$\star\star\star$]\label{exo:b2:search-and-unemployment:11}
The values of employment and unemployment satisfy $rW = w + s(U - W)$ and $rU = z + f(\theta)(W - U)$; those of a filled job and a vacancy satisfy $rJ = p - w - sJ$ and $rV = -c + q(\theta)(J - V)$. With $V = 0$ and
$W - U = \beta(W - U + J)$, derive \eqref{eq:b2:search-and-unemployment:wage}.
\end{exercise}

\begin{exercise}[$\star\star\star$]\label{exo:b2:search-and-unemployment:12}
Show that the elasticity of $\theta$ with respect to $p$ is $\dfrac{p/(p - z)}{\eta + (1 - \eta)\beta f/(r + s + \beta f)}$. Evaluate it for the listing's calibration with $z = 0.4$ and with $z = 0.9$.
\end{exercise}

\section{Problem: Brenland's labour market}\label{sec:b2:search-and-unemployment:problem}

\begin{problem}[{Weekend problem --- Brenland's labour market: flows, matching, the wage and job creation, and a rise in
benefits}]
\label{pb:b2:search-and-unemployment:1}
Per quarter, Brenland's jobs end at rate $s = 0.05$ and the interest rate is $r = 0.01$. Matches are $m = 0.6u^{0.5}v^{0.5}$. A job produces $p = 1$; an open vacancy costs $c = 0.5$; the value of
unemployment is $z = 0.6$; workers' \omterm{def:b2:search-and-unemployment:nash-bargaining}{bargaining power} is $\beta = 0.5$.

\medskip
\noindent\textbf{Part I --- Flows.}
\begin{enumerate}
  \item Write the law of motion of unemployment.
  \item If the job-finding rate were 0.45, what would steady-state unemployment be?
  \item How many quarters, on average, does a spell of unemployment then last?
  \item And a job?
  \item Why does a higher separation rate shift the Beveridge curve outward?
\end{enumerate}

\medskip
\noindent\textbf{Part II --- Matching.}
\begin{enumerate}[resume]
  \item Write $f(\theta)$ and $q(\theta)$.
  \item Show that the \omterm{def:b2:search-and-unemployment:matching-function}{matching function} has constant returns.
  \item What tightness gives a job-finding rate of 0.48?
  \item What is the \omterm{def:b2:search-and-unemployment:vacancy-filling}{vacancy-filling rate} at that tightness?
  \item Write the Beveridge curve as $v$ in terms of $u$.
\end{enumerate}

\medskip
\noindent\textbf{Part III --- The wage and job creation.}
\begin{enumerate}[resume]
  \item Write the \omterm{def:b2:search-and-unemployment:job-creation}{job-creation condition} as $w$ in terms of $\theta$.
  \item Write the wage curve.
  \item Write the equation for $x = \sqrt\theta$ and solve it.
  \item Find the equilibrium tightness, job-finding rate and \omterm{def:b2:search-and-unemployment:vacancy-filling}{vacancy-filling rate}.
  \item Find unemployment and vacancies.
  \item Find the wage, and check it in both curves.
  \item A rise in productivity to 1.02: find tightness and unemployment.
  \item What is the elasticity of tightness with respect to productivity here?
\end{enumerate}

\medskip
\noindent\textbf{Part IV --- A rise in benefits.} The value of unemployment rises to $z = 0.7$.
\begin{enumerate}[resume]
  \item Which curve shifts, and which way?
  \item Find the new tightness and job-finding rate.
  \item Find the new unemployment rate.
  \item Find the new wage. Why does it rise so little?
  \item Why do vacancies fall?
  \item What channel through which benefits affect unemployment does this model leave out, which the studies of the record
  measure?
  \item State the result: equilibrium tightness, unemployment and wage, and the effect of the benefit rise.
\end{enumerate}
\end{problem}
